% !TEX root = ../main.tex
\chapter{Resultados}

Este capítulo reúne la evidencia experimental de los experimentos de
reconstrucción, clasificación de WSI, segmentación de WSI a resolución de
parche y análisis geométrico del espacio latente. Se distingue entre el
conjunto fijo de validación usado para diagnóstico durante el desarrollo, las
pruebas selladas y los análisis exploratorios posteriores. Las diferencias se
expresan como VAE base menos VAE-\(\mathrm{SO}(2)\), salvo indicación contraria.

\section{Configuración experimental final}

Los resultados corresponden a los puntos de control de la actualización
\(60\,000\) de ambos VAE. Los codificadores y decodificadores permanecieron
congelados durante todas las tareas posteriores. La reconstrucción sellada
incluyó \(67\,138\) parches de 23 WSI; la clasificación de diagnóstico usó las
bolsas completas de esas mismas 23 WSI, con \(261\,168\) posiciones; y la
segmentación a resolución de parche evaluó \(31\,572\) posiciones de regiones
anotadas. Cada experimento supervisado tuvo una sola inicialización por rama,
por lo que los intervalos cuantifican remuestreo de WSI, no variación entre
entrenamientos.

La Figura~\ref{fig:entrenamiento-vae} resume las diez épocas de los dos VAE.
Ambas ramas completaron el presupuesto previsto. Las curvas se muestran como
descripción del proceso de optimización; no se usaron para cambiar
retrospectivamente el punto de control evaluado.

\begin{figure}[p]
    \centering
    \includegraphics[width=0.88\textwidth]{vae_training_loss_panels.png}
    \caption{Evolución del objetivo compuesto y del término MAE durante
    \(60\,000\) actualizaciones. Las curvas de entrenamiento están suavizadas y
    los marcadores corresponden a las evaluaciones periódicas de validación.
    Elaboración propia a partir de los registros de entrenamiento.}
    \label{fig:entrenamiento-vae}
\end{figure}

\begin{figure}[p]
    \ContinuedFloat
    \centering
    \includegraphics[width=0.88\textwidth]{vae_training_quality_panels.png}
    \caption[]{Evolución del entrenamiento de los VAE (continuación): pérdida
    estructural \(1-\mathrm{SSIM}\) y PSNR sobre los 25 parches fijos de
    validación. Estos parches se usaron como diagnóstico de desarrollo y no
    como prueba sellada.}
\end{figure}

\begin{figure}[p]
    \ContinuedFloat
    \centering
    \includegraphics[width=0.88\textwidth]{vae_training_control_panels.png}
    \caption[]{Evolución del entrenamiento de los VAE (continuación): razón de
    equivarianza latente de desarrollo y calendario de la tasa de aprendizaje.
    La razón se presenta como diagnóstico de desarrollo, no como criterio para
    escoger un punto de control posterior.}
\end{figure}

\section{Reconstrucción y aprendizaje de representaciones}

\subsection{Inspección cualitativa y conjunto fijo de desarrollo}

Los dos modelos preservaron la organización general del tejido, los núcleos y
las variaciones principales de tinción en los 25 parches fijos. La
Figura~\ref{fig:reconstrucciones-fixed25} permite comparar las mismas entradas
con ambas reconstrucciones. Este panel fue observado repetidamente durante el
desarrollo; por ello, sirve para diagnóstico cualitativo y no como prueba
independiente.

\begin{figure}[htbp]
    \centering
    \includegraphics[width=\textwidth]{vae_reconstructions_fixed25.png}
    \caption{Entradas y reconstrucciones pareadas del conjunto fijo de 25
    parches. Cada columna conserva el mismo parche para el VAE base y el
    VAE-\(\mathrm{SO}(2)\). Elaboración propia a partir de las reconstrucciones
    registradas durante el entrenamiento.}
    \label{fig:reconstrucciones-fixed25}
\end{figure}

El mosaico permite revisar diversidad, pero comprime los detalles celulares.
La Figura~\ref{fig:reconstruccion-ampliada} muestra a resolución mayor el
parche fijo 12, que también se utilizó en los análisis de rotación. La
comparación ampliada ilustra, para este parche, que las diferencias entre modelos son sutiles en
la reconstrucción sin transformar y que ambas ramas suavizan parte de la
textura fina respecto de la entrada.

\begin{figure}[htbp]
    \centering
    \begin{subfigure}[t]{0.30\textwidth}
        \centering
        \includegraphics[width=\linewidth]{reconstruction_zoom_original.png}
        \caption{Parche original.}
    \end{subfigure}\hfill
    \begin{subfigure}[t]{0.30\textwidth}
        \centering
        \includegraphics[width=\linewidth]{reconstruction_zoom_normal.png}
        \caption{VAE base.}
    \end{subfigure}\hfill
    \begin{subfigure}[t]{0.30\textwidth}
        \centering
        \includegraphics[width=\linewidth]{reconstruction_zoom_so2.png}
        \caption{VAE-\(\mathrm{SO}(2)\).}
    \end{subfigure}
    \caption{Ampliación del parche fijo 12 y sus reconstrucciones finales sin
    transformar. Las tres imágenes se muestran a la misma escala y provienen
    de los arreglos numéricos conservados para el análisis de transformación
    latente.}
    \label{fig:reconstruccion-ampliada}
\end{figure}

En este conjunto fijo, el VAE base obtuvo MAE
\(0,0608\pm0,0127\), MSE \(0,0071\pm0,0028\), PSNR
\(27,898\pm1,839\) dB y SSIM \(0,7308\pm0,0637\). El
VAE-\(\mathrm{SO}(2)\) obtuvo, respectivamente,
\(0,0624\pm0,0130\), \(0,0075\pm0,0029\),
\(27,645\pm1,826\) dB y \(0,7170\pm0,0637\). Las dispersiones son
desviaciones estándar poblacionales entre los 25 parches y no intervalos de
incertidumbre.

\subsection{Prueba sellada de reconstrucción}

La Tabla~\ref{tab:reconstruccion-sellada} presenta las métricas sobre los
\(67\,138\) parches de prueba. En las cuatro medidas, el punto estimado favoreció
ligeramente al VAE base. La diferencia primaria de MAE fue \(-0,001358\), con
intervalo bootstrap por WSI del 95\% \([-0,002213;0,000046]\). Como el intervalo
incluyó cero, esta prueba no sustenta una diferencia concluyente en el criterio
primario. El VAE base tuvo menor MAE media en 22 de las 23 WSI, pero ese conteo
descriptivo tampoco incorpora variación entre entrenamientos.

\begin{table}[htbp]
    \centering
    \caption{Reconstrucción en la prueba sellada. Se reporta media
    \(\pm\) desviación estándar entre parches.}
    \label{tab:reconstruccion-sellada}
    \small
    \begin{tabular}{@{}lrrrr@{}}
        \toprule
        Modelo & MAE & MSE & PSNR (dB) & SSIM \\
        \midrule
        VAE base & \(0,0629\pm0,0239\) & \(0,0080\pm0,0074\) & \(27,923\pm2,721\) & \(0,7575\pm0,0842\) \\
        VAE-\(\mathrm{SO}(2)\) & \(0,0642\pm0,0223\) & \(0,0083\pm0,0067\) & \(27,658\pm2,595\) & \(0,7436\pm0,0862\) \\
        \bottomrule
    \end{tabular}
\end{table}

La Figura~\ref{fig:reconstruccion-sellada-parches} muestra por separado las
cuatro distribuciones por parche, mientras la
Figura~\ref{fig:reconstruccion-sellada-wsi} presenta la comparación pareada por
WSI. PSNR y SSIM, definidos como medidas secundarias, también tuvieron mejores
valores descriptivos para el VAE base. En conjunto, se observó una diferencia
pequeña en la medida principal de esta prueba. El experimento no fue diseñado
como prueba de equivalencia o no inferioridad, ni permite atribuir causalmente
la diferencia a la equivarianza porque también cambiaron la parametrización y
el costo de cómputo.

\begin{figure}[p]
    \centering
    \begin{subfigure}[t]{0.96\textwidth}
        \centering
        \includegraphics[width=\linewidth]{vae_test_patch_mae.png}
        \caption{Error absoluto medio en el dominio normalizado.}
    \end{subfigure}

    \begin{subfigure}[t]{0.96\textwidth}
        \centering
        \includegraphics[width=\linewidth]{vae_test_patch_mse.png}
        \caption{Error cuadrático medio en el dominio normalizado.}
    \end{subfigure}
    \caption{Distribuciones de error de reconstrucción por parche sobre las 23
    WSI de prueba. Las cajas describen parches correlacionados dentro de cada
    lámina; no son intervalos de confianza. Elaboración propia a partir de la
    evaluación de reconstrucción en la prueba sellada.}
    \label{fig:reconstruccion-sellada-parches}
\end{figure}

\begin{figure}[p]
    \ContinuedFloat
    \centering
    \begin{subfigure}[t]{0.96\textwidth}
        \centering
        \includegraphics[width=\linewidth]{vae_test_patch_psnr.png}
        \caption{Relación señal a ruido pico por imagen.}
    \end{subfigure}

    \begin{subfigure}[t]{0.96\textwidth}
        \centering
        \includegraphics[width=\linewidth]{vae_test_patch_ssim.png}
        \caption{Índice de similitud estructural por imagen.}
    \end{subfigure}
    \caption[]{Distribuciones de calidad de reconstrucción por parche
    (continuación). Valores mayores indican mejor reconstrucción en PSNR y
    SSIM.}
\end{figure}

\begin{figure}[p]
    \centering
    \includegraphics[width=\textwidth]{vae_test_paired_wsi_mae_readable.png}
    \caption{MAE media pareada de las 23 WSI de prueba. Cada segmento une las
    medias de una misma WSI en el VAE base y el VAE-\(\mathrm{SO}(2)\), por lo
    que no confunde láminas con diferente número de parches. Elaboración propia
    a partir de la evaluación de reconstrucción en la prueba sellada.}
    \label{fig:reconstruccion-sellada-wsi}
\end{figure}

\section{Clasificación a nivel de WSI}

El clasificador MIL local-global se entrenó por separado, con la misma
arquitectura y protocolo, sobre cada representación congelada. La
Figura~\ref{fig:desarrollo-mil} muestra el desarrollo supervisado y el punto
seleccionado en cada rama. Las bolsas se procesaron completas, incluida la más
grande de \(32\,595\) parches.

\begin{figure}[p]
    \centering
    \includegraphics[width=\textwidth]{mil_training_f1.png}
    \caption{Evolución del entrenamiento y validación del clasificador MIL
    local-global: macro-F1 de validación. El punto de control de cada rama se
    eligió por esta métrica, con entropía cruzada como desempate. El azul
    identifica el VAE base y el naranja el VAE-\(\mathrm{SO}(2)\). Elaboración
    propia a partir de los registros del experimento.}
    \label{fig:desarrollo-mil}
\end{figure}

\clearpage

\begin{figure}[p]
    \ContinuedFloat
    \centering
    \includegraphics[width=\textwidth]{mil_training_loss.png}
    \caption[]{Evolución del entrenamiento y validación del clasificador MIL
    local-global (continuación): entropía cruzada. Las líneas continuas
    corresponden a entrenamiento y las discontinuas a validación.}
\end{figure}

\clearpage

En la prueba sellada, el VAE-\(\mathrm{SO}(2)\) tuvo mejores estimaciones
puntuales de macro-F1, exactitud y entropía cruzada; el VAE base tuvo una
exactitud balanceada ligeramente mayor. La Tabla~\ref{tab:mil-sellado} muestra
que el intervalo de la diferencia primaria de macro-F1 atravesó cero. La
entropía cruzada fue menor para la rama equivariante y su intervalo exploratorio
no atravesó cero, pero esto no reemplaza el criterio primario ni corrige la
incertidumbre por una sola inicialización.

\begin{table}[htbp]
    \centering
    \caption{Clasificación de diagnóstico en 23 WSI de prueba sellada. La
    diferencia es VAE base menos VAE-\(\mathrm{SO}(2)\).}
    \label{tab:mil-sellado}
    \small
    \begin{tabular}{@{}lrrrr@{}}
        \toprule
        Métrica & VAE base & VAE-\(\mathrm{SO}(2)\) & Diferencia & IC 95\% \\
        \midrule
        Macro-F1 & 0,389 & 0,483 & \(-0,095\) & \([-0,207;0,050]\) \\
        Exactitud balanceada & 0,510 & 0,503 & \(0,007\) & \([-0,167;0,180]\) \\
        Exactitud & 0,522 & 0,609 & \(-0,087\) & \([-0,217;0,043]\) \\
        Entropía cruzada & 1,323 & 0,993 & \(0,330\) & \([0,135;0,524]\) \\
        \bottomrule
    \end{tabular}
\end{table}

La prueba contenía 5 WSI CC, 6 EC, 8 HGSC, 2 LGSC y 2 MC. La escasa presencia
de LGSC y MC hizo inestables sus medidas por clase. En particular, el VAE base
no recuperó casos EC ni LGSC, mientras que el VAE-\(\mathrm{SO}(2)\) no recuperó
casos MC. Las métricas globales se presentan en la
Figura~\ref{fig:mil-sellado-global}; las matrices de confusión, en la
Figura~\ref{fig:mil-sellado-confusiones}; y los F1 por clase, en la
Figura~\ref{fig:mil-sellado-f1-clases}. El protocolo de inferencia no
almacenó los pesos de atención por parche. Por ello no se reportan mapas de
atención empíricos: recalcularlos con una ejecución posterior no describiría
exactamente la inferencia con la que se obtuvieron estos resultados.

\begin{figure}[p]
    \centering
    \includegraphics[width=\textwidth]{mil_test_metrics.png}
    \caption{Métricas globales de diagnóstico en la prueba sellada de 23 WSI.
    Los valores puntuales deben interpretarse junto con los intervalos por WSI
    de la Tabla~\ref{tab:mil-sellado}. Elaboración propia a partir de las
    predicciones de prueba.}
    \label{fig:mil-sellado-global}
\end{figure}

\begin{figure}[p]
    \centering
    \includegraphics[width=0.94\textwidth]{mil_test_confusion_normal.png}
    \caption{Matrices de confusión en la prueba sellada: clasificador entrenado
    sobre las representaciones del VAE base. Cada celda presenta el conteo y el
    porcentaje dentro de la clase real indicada por la fila. Ambos modelos se
    evaluaron sobre las mismas 23 WSI. Elaboración propia a partir de las
    predicciones de prueba.}
    \label{fig:mil-sellado-confusiones}
\end{figure}

\begin{figure}[p]
    \ContinuedFloat
    \centering
    \includegraphics[width=0.94\textwidth]{mil_test_confusion_so2.png}
    \caption[]{Matrices de confusión en la prueba sellada (continuación):
    clasificador entrenado sobre las representaciones del
    VAE-\(\mathrm{SO}(2)\).}
\end{figure}

\clearpage

\begin{figure}[p]
    \centering
    \includegraphics[width=\textwidth]{mil_test_per_class_f1.png}
    \caption{F1 por subtipo histológico en la prueba sellada. El soporte
    pequeño de LGSC y MC impide interpretar una barra aislada como una
    estimación estable de desempeño poblacional. Elaboración propia a partir de
    las predicciones de prueba.}
    \label{fig:mil-sellado-f1-clases}
\end{figure}

\clearpage

\section{Segmentación de WSI a resolución de parche}

La tarea de tejido asignó tumor, estroma o necrosis a cada celda seleccionada
de la cuadrícula de parches. Es, por tanto, una segmentación de WSI a resolución
de \(256\times256\) píxeles, pero no exhaustiva: solo se evaluaron regiones con
anotación de alta pureza. La prueba contenía \(21\,796\) parches de tumor,
\(8\,500\) de estroma y \(1\,276\) de necrosis. Esta última clase apareció en
solo cinco WSI.

La Tabla~\ref{tab:eficiencia-etiquetas} resume el macro-F1 para los cinco
presupuestos anidados. La rama equivariante tuvo una estimación mayor entre 250
y \(2\,500\) etiquetas por clase, mientras la rama base fue mayor con todas las
\(5\,671\) disponibles. Solo el intervalo simultáneo del presupuesto 500
excluyó cero. El área normalizada bajo la curva de eficiencia fue 0,615 para el
VAE base y 0,631 para el VAE-\(\mathrm{SO}(2)\), con diferencia
\(-0,016\;[-0,073;0,040]\); por consiguiente, la familia completa no demostró
una ventaja concluyente.

\begin{table}[htbp]
    \centering
    \caption{Macro-F1 en la segmentación dispersa de WSI a resolución de
    parche. Los intervalos simultáneos corresponden a la diferencia VAE base
    menos VAE-\(\mathrm{SO}(2)\).}
    \label{tab:eficiencia-etiquetas}
    \small
    \begin{tabular}{@{}rrrrr@{}}
        \toprule
        Etiquetas/clase & VAE base & VAE-\(\mathrm{SO}(2)\) & Diferencia & IC simultáneo 95\% \\
        \midrule
        250 & 0,508 & 0,523 & \(-0,015\) & \([-0,072;0,041]\) \\
        500 & 0,511 & 0,575 & \(-0,064\) & \([-0,121;-0,008]\) \\
        \(1\,000\) & 0,582 & 0,598 & \(-0,016\) & \([-0,072;0,041]\) \\
        \(2\,500\) & 0,701 & 0,714 & \(-0,013\) & \([-0,069;0,044]\) \\
        \(5\,671\) & 0,766 & 0,710 & \(0,055\) & \([-0,001;0,112]\) \\
        \bottomrule
    \end{tabular}
\end{table}

La Figura~\ref{fig:tejido-resultados} evidencia que la relación entre
representación y cantidad de etiquetas no fue monótona. También muestra la
dependencia por clase: necrosis tuvo el menor soporte y la mayor sensibilidad a
las cinco WSI que la contenían. Los resultados describen únicamente las
regiones anotadas y no asignan una clase a los sectores negros o sin máscara.

\clearpage

\begin{figure}[p]
    \centering
    \includegraphics[width=\textwidth]{tissue_training_f1.png}
    \caption{Mejor macro-F1 de validación alcanzado por el clasificador de
    tejido en cada presupuesto de etiquetas. El azul identifica el VAE base y
    el naranja el VAE-\(\mathrm{SO}(2)\). La curva pertenece al desarrollo y se
    presenta separada de los resultados de prueba sellada. Elaboración propia a
    partir de los registros del experimento.}
    \label{fig:tejido-desarrollo}
\end{figure}

\clearpage

\begin{figure}[p]
    \centering
    \includegraphics[width=\textwidth]{tissue_label_efficiency.png}
    \caption{Prueba sellada de segmentación a resolución de parche: macro-F1
    frente al presupuesto por clase. El eje de etiquetas refleja subconjuntos
    balanceados, anidados y estratificados por WSI. Elaboración propia a partir
    de las predicciones de prueba.}
    \label{fig:tejido-resultados}
\end{figure}

\clearpage

\begin{figure}[p]
    \ContinuedFloat
    \centering
    \includegraphics[width=\textwidth]{tissue_test_per_class.png}
    \caption[]{Prueba sellada de segmentación a resolución de parche
    (continuación): desempeño de tumor, estroma y necrosis en los cinco
    presupuestos de etiquetas.}
\end{figure}

\clearpage

\section{Análisis del espacio latente}

\subsection{Transformaciones prescritas y reconstrucciones}

La primera comprobación aplicó giros exactos de 90, 180 y 270 grados al mapa
latente espacial y decodificó el resultado. La Figura~\ref{fig:rotacion-latente}
muestra que el VAE-\(\mathrm{SO}(2)\) siguió la acción prescrita en estos cuartos
de vuelta, mientras que aplicar la misma operación geométrica al latente del
VAE base produjo cambios y artefactos no controlados. Esta es una comprobación
cualitativa sobre parches fijos, no una comparación poblacional de calidad.

\begin{figure}[htbp]
    \centering
    \includegraphics[width=0.94\textwidth]{vae_rotated_input_vs_latent.png}
    \caption{Reconstrucciones al rotar la entrada y al transformar el latente.
    Los cuartos de vuelta son controles exactos de la acción espacial. En el
    VAE base, rotar el mapa latente es una intervención para la cual no fue
    entrenado;
    en el VAE-\(\mathrm{SO}(2)\), corresponde a la acción prescrita. Fuente:
    análisis de transformación latente.}
    \label{fig:rotacion-latente}
\end{figure}

La ampliación de la Figura~\ref{fig:rotacion-latente-ampliada} aísla el giro de
90 grados y las dos intervenciones latentes. En el VAE base aparecen mezclas de
textura y contornos que no corresponden al objetivo rotado; en el
VAE-\(\mathrm{SO}(2)\), la transformación espacial del posterior conserva la
organización de la reconstrucción. Este ejemplo hace visible el mecanismo que
el mosaico completo resume, pero no reemplaza las medidas sobre los 25 parches.

\begin{figure}[htbp]
    \centering
    \begin{subfigure}[t]{0.30\textwidth}
        \centering
        \includegraphics[width=\linewidth]{latent_action_zoom_target.png}
        \caption{Objetivo a 90 grados.}
    \end{subfigure}\hfill
    \begin{subfigure}[t]{0.30\textwidth}
        \centering
        \includegraphics[width=\linewidth]{latent_action_zoom_normal.png}
        \caption{Acción sobre el latente base.}
    \end{subfigure}\hfill
    \begin{subfigure}[t]{0.30\textwidth}
        \centering
        \includegraphics[width=\linewidth]{latent_action_zoom_so2.png}
        \caption{Acción sobre el latente \(\mathrm{SO}(2)\).}
    \end{subfigure}
    \caption{Ampliación del parche fijo 12 al transformar el latente 90 grados
    y decodificarlo. Las imágenes provienen de los arreglos numéricos del
    análisis de transformación latente.}
    \label{fig:rotacion-latente-ampliada}
\end{figure}

La evaluación cuantitativa distinguió entre transformar el latente y
canonizarlo antes de decodificar. Para rotaciones interpoladas de cinco grados,
la mediana de la razón de error de acción fue 0,713 en el VAE base y 0,550 en
el VAE-\(\mathrm{SO}(2)\); la razón de canonización fue 0,661 y 0,513,
respectivamente. La rama equivariante fue menor en los 25 parches, pero no
alcanzó el criterio absoluto predefinido de 0,5. En cambio, para el control
exacto de cuartos de vuelta, la razón de acción descendió de 0,709453 en el VAE
base a \(2,18\times10^{-6}\) en el VAE-\(\mathrm{SO}(2)\), y la de
canonización de 0,710638 a 0,088854.

\begin{figure}[htbp]
    \centering
    \begin{subfigure}[t]{\textwidth}
        \centering
        \includegraphics[width=0.82\linewidth]{latent_population_ratios.png}
        \caption{Razones poblacionales para rotaciones interpoladas.}
    \end{subfigure}

    \begin{subfigure}[t]{\textwidth}
        \centering
        \includegraphics[width=0.96\linewidth]{latent_exact_d4.png}
        \caption{Control exacto a 90, 180 y 270 grados.}
    \end{subfigure}
    \caption{Respuesta de las reconstrucciones a transformaciones del latente.
    Valores menores indican mejor acuerdo con la acción o canonización
    prescrita. El control exacto verifica la compatibilidad del decodificador
    equivariante con \(C_4\), no por sí solo equivarianza continua de todo el
    autoencodificador. Fuente: análisis de transformación latente.}
    \label{fig:razones-latentes}
\end{figure}

La distinción anterior resultó esencial. Al evaluar la composición completa
\(D(\boldsymbol{\mu}(R_\theta x))\) frente a
\(R_\theta D(\boldsymbol{\mu}(x))\), que también incluye al codificador y la
interpolación de la entrada, la mediana fue 0,174 para el VAE base y 0,210 para
el VAE-\(\mathrm{SO}(2)\); el VAE base fue menor en los 25 parches. Por ello,
el éxito del control \(C_4\) no se extrapola a una ventaja de extremo a extremo
para rotaciones continuas.

\subsection{Geometría local y organización espacial}

En órbitas muestreadas cada grado, la medida corregida de geometría local fue
0,649 para el VAE base y 0,601 para el VAE-\(\mathrm{SO}(2)\), una reducción
relativa de 7,38\%. Aunque los 25 parches favorecieron a la rama equivariante a
esa resolución, el efecto no alcanzó el margen predefinido de 10\% y la
dirección no fue estable al usar pasos de cinco grados: el VAE-\(\mathrm{SO}(2)\)
fue favorable en solo 8 de 25 parches. Tampoco se obtuvieron resultados que
demostraran una acción reducida compartida o una factorización global del
latente.

La Figura~\ref{fig:orbita-latente-parche} muestra cómo se obtuvo esa evidencia
para el parche fijo 12. Cada punto representa el posterior completo dentro del
disco central para un ángulo, proyectado sobre un PCA ajustado con las 360
vistas del mismo modelo y parche. La trayectoria del VAE-\(\mathrm{SO}(2)\) es
visualmente más regular en este ejemplo; la curva inferior recuerda que la
forma de la proyección no sustituye el residual directo ni permite comparar
las escalas de dos PCA independientes.

\begin{figure}[htbp]
    \centering
    \includegraphics[width=0.96\textwidth]{latent_orbit_single_patch.png}
    \caption{Órbita latente de un parche al rotar la entrada entre 0 y 359
    grados. Arriba se muestran los dos PCA independientes; abajo, el residual
    directo para el parche y el control exacto agregado de cuartos de vuelta.
    Fuente: análisis de órbitas latentes.}
    \label{fig:orbita-latente-parche}
\end{figure}

La proyección PCA espacial de la Figura~\ref{fig:pca-latente-espacial} mostró
estructura local en ambos modelos. La RMS relativa de aristas tuvo mediana
1,4567 para el VAE base y 1,4676 para el VAE-\(\mathrm{SO}(2)\); por tanto, esta
medida no mostró mayor suavidad espacial de la rama equivariante. Un sondeo RGB
puntual explicó \(R^2=0,286\) y \(R^2=0,189\), respectivamente. Como los PCA se
ajustaron por modelo, los colores no alinean canales entre redes.

\begin{figure}[htbp]
    \centering
    \includegraphics[width=0.70\textwidth]{latent_spatial_pca_maps.png}
    \caption{Proyección PCA de los mapas latentes espaciales para cuatro
    parches fijos. Cada color representa los tres primeros componentes bajo un
    ajuste común por modelo; se usa para visualizar organización interna y no
    como correspondencia de canales entre redes. Fuente: análisis de órbitas y
    estructura latente.}
    \label{fig:pca-latente-espacial}
\end{figure}

\begin{figure}[htbp]
    \ContinuedFloat
    \centering
    \begin{subfigure}[t]{\textwidth}
        \centering
        \includegraphics[width=0.96\linewidth]{latent_spatial_pca_paired.png}
        \caption{RMS relativa de aristas, pareada sobre los 25 parches.}
    \end{subfigure}

    \begin{subfigure}[t]{\textwidth}
        \centering
        \includegraphics[width=0.96\linewidth]{latent_spatial_pca_delta.png}
        \caption{Diferencia VAE-\(\mathrm{SO}(2)\) menos VAE base por parche.}
    \end{subfigure}
    \caption[]{Proyección PCA espacial (continuación): medidas cuantitativas de
    rugosidad local. Un valor menor no equivale por sí solo a una mejor
    representación, pues también puede resultar de suavizado o colapso.}
\end{figure}

En síntesis, los resultados latentes separan una propiedad verificada de una
hipótesis aún abierta. La compatibilidad numérica del decodificador
VAE-\(\mathrm{SO}(2)\) con la acción discreta \(C_4\) en cuartos de vuelta
quedó verificada, y a un grado se
observó una mejora local descriptiva. No quedó demostrada una organización
global más simple, una ventaja continua de extremo a extremo ni una mejora
uniforme de las tareas posteriores. El análisis funcional planificado deberá
evaluarse antes de formular conclusiones adicionales.

\section{Síntesis de resultados por objetivo}

La Tabla~\ref{tab:sintesis-objetivos} relaciona los objetivos específicos con
la evidencia disponible y evita que los resultados exploratorios sustituyan
los criterios principales.

\begin{table}[htbp]
    \centering
    \caption{Síntesis de evidencia por objetivo específico.}
    \label{tab:sintesis-objetivos}
    \small
    \begin{tabularx}{\textwidth}{@{}p{0.20\textwidth}XX@{}}
        \toprule
        Objetivo & Evidencia y resultado principal & Limitación principal \\
        \midrule
        Selección de datos & UBC-OCEAN permitió entrenamiento no supervisado,
        diagnóstico de cinco subtipos y evaluación de tres tejidos. & Las
        máscaras son suplementarias y no cubren exhaustivamente las WSI. \\
        Formulación e implementación & Se construyeron VAE residuales pareados,
        uno con convoluciones restringidas a \(F_0/F_1\), y clasificadores
        idénticos sobre sus latentes congelados. & La rama equivariante tiene
        menos parámetros, pero mayor costo de convolución; capacidad y simetría
        no quedaron aisladas causalmente. \\
        Evaluación & La reconstrucción favoreció levemente al VAE base; WSI y
        algunos presupuestos bajos favorecieron puntualmente al
        VAE-\(\mathrm{SO}(2)\); se verificó su control exacto \(C_4\). & Una
        inicialización por rama, 23 WSI de prueba, clases minoritarias escasas y
        análisis funcional latente pendiente. \\
        \bottomrule
    \end{tabularx}
\end{table}
